\chapter{威尔逊费米子 (WF)}

在第~8 节中已经说明，格点作用量不是唯一的。人们可以自由地在作用量中添加任意数量的无关（irrelevant）算符，因为这些算符不改变连续极限。Wilson 对加倍问题的解决方案是添加一个维数为五的算符 $a r \bar\psi \square \psi$，由此在 $p_\mu = \pi$ 处的额外十五个费米子种类获得与 $r/a$ 成正比的质量 [56]。威尔逊作用量为
\begin{align}
A^W &= m_q \sum_x \bar\psi(x)\psi(x) \nonumber\\
&\quad + \frac{1}{2a}\sum_{x,\mu} \bar\psi(x)\gamma_\mu \bigl[ U_\mu(x)\psi(x+\hat{\mu}) - U_\mu^\dagger(x-\hat{\mu})\psi(x-\hat{\mu}) \bigr] \nonumber\\
&\quad - \frac{r}{2a}\sum_{x,\mu} \bar\psi(x) \bigl[ U_\mu(x)\psi(x+\hat{\mu}) - 2\psi(x) + U_\mu^\dagger(x-\hat{\mu})\psi(x-\hat{\mu}) \bigr] \nonumber\\
&= \frac{(m_q a + 4r)}{a}\sum_x \bar\psi(x)\psi(x) \nonumber\\
&\quad + \frac{1}{2a}\sum_x \bar\psi \bigl[ (\gamma_\mu - r) U_\mu(x)\psi(x+\hat{\mu}) - (\gamma_\mu + r) U_\mu^\dagger(x-\hat{\mu})\psi(x-\hat{\mu}) \bigr] \nonumber\\
&\equiv \sum_{x,y} \bar\psi^L_x M^W_{xy} \psi^L_y
\end{align}
其中相互作用矩阵 $M^W$ 通常写成
\begin{equation}
M^W_{x,y}[U]\,a = \delta_{xy} - \kappa \sum_\mu \bigl[ (r - \gamma_\mu) U_{x,\mu}\delta_{x,y-\mu} + (r + \gamma_\mu) U^\dagger_{x-\mu,\mu}\delta_{x,y+\mu} \bigr]
\end{equation}
经过重新标度
\begin{align}
\kappa &= 1/(2m_q a + 8r) \nonumber\\
\psi^L &= \sqrt{m_q a + 4r}\ \psi = \psi/\sqrt{2\kappa} \,.
\end{align}

在这种形式下，威尔逊作用量可以解释为一种统计力学模型。局域项的作用类似于惯性，将费米子保持在同一个格点上，而非局域项则使费米子以强度 $\kappa$ 跳到最近邻格点。在跳跃过程中，费米子在自旋空间获得 $(\gamma_\mu - r)$ 的扭转，在色空间获得 $U_\mu$ 的扭转。$r = 1$ 的情形是特殊的，因为自旋扭转 $1 \pm \gamma_\mu$ 是一个秩为 2 的投影算符，
\begin{align}
\Bigl( \frac{1 \pm \gamma_\mu}{2} \Bigr)^2 &= \frac{1 \pm \gamma_\mu}{2} \nonumber\\
\mathrm{Tr} \Bigl( \frac{1 \pm \gamma_\mu}{2} \Bigr) &= 2 \,.
\end{align}
稍后我们将看到，对于 $r = 1$，色散关系只有一个分支，即自由场极限下没有加倍态。

由式~10.3 可以注意到，对于自由理论，夸克质量用格点参数 $\kappa$ 和 $r$ 表示为
\begin{equation}
m_q a = \frac{1}{2\kappa} - 4r \equiv \frac{1}{2\kappa} - \frac{1}{2\kappa_c}
\end{equation}
在 $\kappa = \kappa_c \equiv 1/8r$ 处为零。对于相互作用理论（$U_\mu(x) \neq 1$），我们将继续定义 $m_q a = 1/2\kappa - 1/2\kappa_c$，但附带条件是 $\kappa_c$ 依赖于 $a$。$\kappa_c$ 与 $\kappa$ 无关的重整化意味着夸克质量同时具有乘法重整化和加法重整化。这是由正比于 $r$ 的无关项显式破坏手征对称性所造成的后果。Wilson 对加倍态的解决方案付出了高昂的代价——手征对称性在 $O(a)$ 量级上的硬破坏。

威尔逊费米子在动量空间中的自由传播子为
\begin{equation}
S_F(p) = M_W^{-1}(p) = \frac{a}{\,1 - 2\kappa \sum_\mu \bigl( r \cos p_\mu a - i\gamma_\mu \sin p_\mu a \bigr)}
\end{equation}
由于维数为 5 的算符（正比于 $r$ 的项），在 $p_\mu a = \pi$ 处的 15 个额外态获得量级为 $2rn/a$ 的质量，其中 $n$ 是 $p_\mu a = \pi$ 的分量个数。

算符 $M$ 满足下列关系：
\begin{align}
\gamma_5 M_W^\dagger \gamma_5 &= M_W, \nonumber\\
\gamma_5 S_F^\dagger (x,y) \gamma_5 &= S_F (y,x) \ , \nonumber\\
M_W^\dagger (\kappa, r) &= M_W (-\kappa, -r) \ .
\end{align}

前两个式子表明 $M^N$ 的「厄米性」性质得以保持。第二个式子将 $x \to y$ 的夸克传播子与 $y \to x$ 的反夸克传播子联系起来。这个重要的恒等式称为厄米性或 $\gamma_5$ 不变性，如第~17.1 节所讨论的，它使得数值计算得到显著简化。$S_F^\dagger$ 中的伴随是相对于每个格点上的自旋和色指标而言的。最后一个关系表明，由于 Wilson 的 $r$ 项，$M^W$ 不是反厄米的。

对朴素费米子作用量所做的与式~7.9 类似的分析表明，由式~10.6 导出的极点质量不同于裸质量，其表达式为
\begin{equation}
m_q^{\mathrm{pole}} a = r(\cosh Ea - 1) + \sinh Ea \,.
\end{equation}
这表明，正如预期的那样，谱量中的离散化修正出现在 $O(a)$ 量级。

\section{威尔逊费米子的性质}

下面简要总结威尔逊费米子的重要性质。

\begin{itemize}
\item 加倍态被赋予一个重的质量 $2 r n /a$，并在连续极限中退耦。

\item 手征对称性被显式破坏。利用式~7.11 所给变换下的不变性推导轴矢沃德恒等式，其一般形式为
\begin{equation}
\Bigl\langle \frac{\partial \delta\mathcal{S}}{\partial \theta}\,\mathcal{O} \Bigr\rangle = \Bigl\langle \frac{\partial \delta\mathcal{O}}{\partial \theta} \Bigr\rangle \,.
\end{equation}
对于 WF，作用量在轴矢变换下的变分为
\begin{equation}
\frac{\partial \delta\mathcal{S}}{\partial \theta} = \partial_\mu A_\mu - 2 m P + r a X,
\end{equation}
其中 $X$ 是来自 Wilson 的 $r$ 项变分的附加项 [57]。因此，一般而言，所有基于轴矢 WI 的关系都将含有 $ra$ 量级的修正，并涉及与那些通常会因手征对称性而消失的算符的混合。

\item 对矩阵元 $\langle 0| T[A_\mu V_\nu V_\rho] |0\rangle$ 的计算表明，ABJ 反常在连续极限中被正确地重现 [58]。Karsten 和 Smit 通过显式的 1 圈计算证明，虽然额外的 15 个态各自贡献正比于 $r$ 的项，但全部十六个态的总贡献与 $r$ 无关，并且等于 ABJ 反常。

\item 夸克质量同时获得加法重整化和乘法重整化。

\item 夸克质量的零点由式~10.5 中的 $\kappa_c$ 决定。在任意给定的 $a$ 处，有两种计算 $\kappa_c$ 的方法。(i) 假定手征关系 $M_\pi^2 \propto m_q$，将$\pi$介子质量计算为 $1/2\kappa$ 的函数，并外推到零。$\pi$介子变为无质量时所对应的 $\kappa$ 值，按定义即为 $\kappa_c$。(ii) 通过比值 $\langle \partial_\mu A_\mu(x) P(0)\rangle / \langle P(x) P(0)\rangle$（基于轴矢沃德恒等式）将夸克质量计算为 $1/2\kappa$ 的函数，并外推到零。这两种估计可能相差 $O(a)$ 量级的修正。这两种计算都需要在规范组态上进行统计平均，并在 $1/2\kappa$ 上进行外推。因此，在单个组态上，矩阵 $M$ 的零模不会正好出现在 $\kappa_c$ 处，而是在 $\kappa_c$ 附近有一个分布。

在全理论中，$\det M$ 抑制具有零模的组态，然而，这种保护在 QQCD 中会丧失。尽管如此，在某些组态（称为「例外」组态）上于 $\kappa < \kappa_c$ 时出现零模的问题，在两种理论中都需要加以解决。

在这些零模附近，即使处于物理区域 $\kappa < \kappa_c$，夸克传播子也会变得奇异。因此，在足够大的统计样本中，人们会遇到「例外」组态，这些组态上的强子关联函数与其余组态相比会显示出数量级上的涨落（取决于夸克质量有多小）。由于这种病态是类 Wilson 作用量的特征，而不是规范组态本身的特征，因此不能简单地将这些组态从统计样本中剔除。此类零模的问题随夸克质量减小和 $\beta$ 减小而加剧。迄今为止使用的两种实用方法是：(i) 在 $m_q > m_q^{*}$ 处模拟，其中 $m_q^{*}$ 足够大，使得零模位置的弥散不会在关联函数中引起大的涨落；(ii) 对于固定的 $m_q$，增大格点体积，因为这降低了遇到例外组态的概率。如果人们感兴趣的是在手征极限附近模拟 QCD，即物理轻夸克，这些补救措施没有帮助。为此需要新的思路。关于此问题的更详细分析以及可能的解决方案见 [59,60]。

\item 自旋和味道自由度与连续统中的 Dirac 费米子一一对应。因此，插值场算符的构造是直接的，例如，$\bar\psi \gamma_5 \psi$ 和 $\bar \psi \gamma_i \psi$ 分别是赝标量介子和矢量介子的插值算符，正如在连续统中一样（见第~17 节）。

\item Wilson 项将离散化误差改变为 $O(a)$ 量级。

\item \textbf{费曼规则}以及方格规范作用量和威尔逊费米子的规范固定细节见 [36]。

\item \textbf{$S_F$ 的对称性}：传播子在分立对称性下的变换与朴素费米子相同。对于给定的背景组态 $U$，这些变换为
\begin{align}
\mathcal{P}: \qquad S_F(x,y,[U]) &\to \gamma_4\,S_F(x^P,y^P,[U^P])\,\gamma_4 \nonumber\\
\mathcal{T}: \qquad S_F(x,y,[U]) &\to \gamma_4\gamma_5\,S_F(x^T,y^T,[U^T])\,\gamma_5\gamma_4 \nonumber\\
\mathcal{C}: \qquad S_F(x,y,[U]) &\to \gamma_4\gamma_2\,S_F^{\mathrm{T}}(x,y,[U^C])\,\gamma_2\gamma_4 \nonumber\\
\mathcal{H}: \qquad S_F(x,y,[U]) &\to \gamma_5\,S_F^\dagger(y,x,[U])\,\gamma_5 \nonumber\\
\mathcal{CPH}: \quad S_F(x,y,[U]) &\to \mathcal{C}\gamma_4\,S_F^\dagger(y^P,x^P,[U^C])\,\gamma_4\mathcal{C}^{-1}
\end{align}
这些关系在确定关联函数的性质（实部与虚部、关于 $p$ 的偶奇性等）时非常有用，如第~17 节所讨论的。

\item 色散关系式~10.8 有以下两个解
\begin{equation}
e^{Ea} = \frac{(ma + r) \pm \sqrt{1 + 2mar + m^2 a^2}}{1+r}
\end{equation}
对于 $r = 0$ 和小 $m$（朴素费米子作用量），这两个解为
\begin{align}
E &= m + O(m^2 a^2) \nonumber\\
E &= -m + \frac{i\pi}{a} + O(m^2 a^2) \,,
\end{align}
其中第二个解对应于欧几里得传播子在 $p_4 = \pi$ 处的极点，即时间加倍态。对于 $r = 1$，正如预期，这个加倍态被消除，唯一的根 $Ea = \log(1 + ma)$ 与物理粒子相关联。对于一般的 $r \neq 1$，第二个根为
\begin{equation}
E = -m + \frac{1}{a} \Bigl( i\pi + \log \frac{1-r}{1+r} \Bigr) + \ldots
\end{equation}
这些不与物理态或加倍态相关联的分支称为鬼分支。我采用一个非常严格的加倍态定义——即在 $m = \vec{p} = 0$ 时满足 $|E| = 0$ 的解。同样由式~10.14 可以清楚地看出，对于 $r = 1$，这个鬼分支被推到无穷远。一般而言，如果作用量耦合了 $(n+1)$ 个时间片，那么色散关系就有 $n$ 个分支。这些根可以是实的、复共轭对，并且/或者具有如上例所示的虚部 $i\pi/a$。这样的附加分支对于改进作用量可能是一个麻烦，例如第~13.2 节讨论的 $D234$ 作用量。目标是得到无加倍态、无鬼分支的理论。这可以通过调整参数（例如在 Wilson 情形中选择 $r = 1$）和/或采用 Lepage 及其合作者所提倡的各向异性格点 $a_t \ll a_s$ [107] 相结合来实现。

\item \textbf{守恒矢量流}：威尔逊作用量在式~7.10 所定义的全局变换下不变。为了推导相应的守恒流（对于简并质量），我们使用标准技巧：计算作用量在依赖于时空的 $\theta(x)$ 变换下的变分，并将 $\delta S$ 中 $\theta(x)$ 的系数与 $\partial_\mu J_\mu$ 联系起来。作用量在这样的无穷小变换下的改变为
\begin{align*}
\delta S &= \kappa \sum_{x,\mu} \bar\psi(x)(\gamma_\mu - r) U_\mu(x) \psi(x + \hat{\mu}) \exp\bigl( i\theta(x) - i\theta(x + \mu) \bigr) \\
&\quad - \kappa \sum_{x,\mu} \bar\psi(x + \mu)(\gamma_\mu + r) U_\mu^\dagger (x) \psi(x) \exp\bigl( i\theta(x + \mu) - i\theta(x) \bigr)
\end{align*}
到 $\theta$ 的一阶，即为
\[
-\sum_{x,\mu} \Bigl[ \bar\psi(x)(\gamma_\mu - r) U_\mu(x) \psi(x + \hat{\mu}) + \bar\psi(x + \mu)(\gamma_\mu + r) U_\mu^\dagger (x) \psi(x) \Bigr] \Bigl[ i \frac{\partial \theta}{\partial x_\mu} \Bigr] .
\]
经分部积分后得到的守恒流为
\begin{equation}
V^{c}_\mu = \bar\psi(x) (\gamma_\mu - r) U_\mu (x) \psi (x + \hat{\mu}) + \bar\psi(x + \mu)(\gamma_\mu + r) U_\mu^\dagger (x) \psi (x) \,.
\end{equation}
$V^{c}_\mu$ 是厄米的，并且在 $r = 0$ 时约化为 1-link 矢量流的对称化形式。在许多应用（如衰变常数）中，人们使用如下定义的局域（味道）流
\begin{align}
2V_\mu(x) &= \bar\psi(x) \gamma_\mu \frac{\lambda^a}{2} U_\mu (x) \psi (x + \hat{\mu}) + \mathrm{h.c.} \nonumber\\
2A_\mu(x) &= \bar\psi(x) \gamma_5\gamma_\mu \frac{\lambda^a}{2} U_\mu (x) \psi (x + \hat{\mu}) + \mathrm{h.c.} \,,
\end{align}
这些流不守恒，因此具有与之关联的非平凡重整化因子 $Z_V$ 和 $Z_A$，在计算矩阵元时必须加以考虑。

\item \textbf{手征凝聚}：如 [57] 中所解释的，手征凝聚 $\langle \bar\psi \psi \rangle$ 与威尔逊夸克传播子的迹 $\langle S_F(0,0) \rangle$ 并不简单地相关。$r$ 项对手征对称性的破坏引入了接触项，这些接触项需要以非微扰的方式从 $S_F(0,0)$ 中减去。这一做法已被证明并不实用。首选的方法是使用连续统中的沃德恒等式
\begin{equation}
\langle \bar\psi\psi \rangle^{\mathrm{WI}} \equiv \langle 0| S_F(0,0) |0\rangle = \lim_{m_q \to 0} m_q \int d^4x\, \langle 0| P(x)P(0) |0\rangle,
\end{equation}
或者 Gell-Mann、Oakes、Renner（GMOR）关系的格点版本 [61]
\begin{equation}
\langle \bar\psi\psi \rangle^{\mathrm{GMOR}} = \lim_{m_q \to 0} - \frac{f_\pi^2 M_\pi^2}{4 m_q} \,.
\end{equation}
两种方法效果的比较见 [62] 的讨论。

\item \textbf{算符混合}：一般而言，在 Callen Symanzik 方程中，不同手性的算符之间存在混合。即使在 $\kappa_c$ 处也是如此，这是手征对称性被显式破坏的结果。这种算符混合给弱相互作用哈密顿量矩阵元的计算带来了严重问题，即沃德恒等式的格点分析以及获得具有正确手征行为的振幅更加复杂。例如，树级 4 费米子算符（与提取 $B_K$、$B_7$ 和 $B_8$ 相关）的矩阵元的手征行为与连续统理论中并不相同。其原因在于，与自发破缺的 $SU_A(n_f)$ 相联系的沃德恒等式在格点理论中会收到 $O(a)$ 量级的修正。Martinelli 教授将更详细地讨论这个问题。
\end{itemize}

在第~13 节中，我将讨论 $WF$ 中的 $O(a)$ 误差如何能以很少的额外努力被降低到 $O(a^2)$ 或更高阶。暂且让我利用威尔逊费米子的这一缺点来引出交错费米子的构造。该构造具有 $O(a^2)$ 量级的误差，保留了足够的手征对称性，从而为涉及 Goldstone $\pi$介子的振幅给出正确的手征行为，但其代价是味道数目增加到原来的 4 倍。
